% !TEX root = ../main.tex
\begin{ouverture}[molecules]{11}{Structure des molécules organiques}{Partie 1 · Constitution et transformations de la matière}
{Squelettes carbonés\enspace·\enspace Familles fonctionnelles\enspace·\enspace Nomenclature\enspace·\enspace Spectroscopie infrarouge et polymères}
\objectif{Passer d'une représentation d'une molécule à une autre ; reconnaître des isomères de constitution.}
\objectif{Identifier les familles fonctionnelles : alcool, aldéhyde, cétone, acide carboxylique, ester, amine, amide, halogénoalcane.}
\objectif{Nommer des molécules simples et retrouver une formule à partir d'un nom.}
\objectif{Exploiter un spectre infrarouge pour identifier des groupes caractéristiques.}
\objectif{Identifier le motif d'un polymère et relier sa masse molaire au degré de polymérisation.}
\prerequis{Formules brutes et semi-développées, liaison covalente (classe de seconde).}
\prerequis{Groupes hydroxyle et carbonyle, spectroscopie infrarouge (classe de première).}
\prerequis{Quantité de matière et masse molaire.}
\end{ouverture}
\essentiel
\essSection{1}{Squelettes carbonés et isomères}
\begin{definition}[Représentations d'une molécule]
La \textbf{formule brute} indique la nature et le nombre des atomes (\ce{C4H10O}) ; la \textbf{formule semi-développée}
explicite les liaisons entre groupes d'atomes ; la \textbf{formule topologique} ne représente que les liaisons du
squelette carboné, chaque extrémité ou brisure de ligne figurant un atome de carbone portant les hydrogènes
nécessaires. Les hétéroatomes (O, N, Cl\,\dots) et leurs hydrogènes sont toujours écrits.
\end{definition}
\schema{11-topologique}
\begin{definition}[Isomères de constitution]
Des \textbf{isomères de constitution} ont la même formule brute mais des enchaînements d'atomes différents : squelettes
différents (chaîne linéaire, ramifiée ou cyclique), positions différentes d'un groupe, ou familles différentes. Leurs
propriétés physiques et chimiques diffèrent.
\end{definition}
\essSection{2}{Familles fonctionnelles et nomenclature}
\begin{savoirfaire}[Nommer une molécule]
Repérer la famille, puis la chaîne carbonée la plus longue qui porte le groupe caractéristique : elle donne le nom de
base (méth-, éth-, prop-, but-, pent-, hex-\,\dots). Numéroter la chaîne pour que le groupe caractéristique ait le plus
petit indice possible, puis indiquer les ramifications (méthyl-, éthyl-\,\dots) avec leurs positions, dans l'ordre
alphabétique.
\end{savoirfaire}
\begin{exemple}
\chemfig[atom sep=2em]{CH_3-CH(-[::60]CH_3)-CH_2-[:0](=[:60]O)-[:-60]OH}\enspace : la chaîne la plus longue portant le
groupe carboxyle compte quatre carbones, numérotés depuis lui ; un groupe méthyle est porté par le carbone 3 :
c'est l'acide 3-méthylbutanoïque.
\end{exemple}
\begin{propriete}[Groupes caractéristiques et familles]
\par\smallskip
{\small\renewcommand{\arraystretch}{1.5}
\begin{tabularx}{\linewidth}{@{}lllY@{}}\toprule
\otTete{Famille} & \otTete{Groupe} & \otTete{Nom} & \otTete{Exemple}\\\midrule
alcool & \chemfig[atom sep=1.6em]{-[:30]OH} \;hydroxyle & \dots-ol & propan-2-ol, \ce{CH3-CHOH-CH3}\\
aldéhyde & \chemfig[atom sep=1.6em]{-[:30](=[:90]O)-[:-30]H} \;carbonyle & \dots-al & éthanal, \ce{CH3-CHO}\\
cétone & \chemfig[atom sep=1.6em]{-[:30](=[:90]O)-[:-30]} \;carbonyle & \dots-one & propanone, \ce{CH3-CO-CH3}\\
acide carboxylique & \chemfig[atom sep=1.6em]{-[:30](=[:90]O)-[:-30]OH} \;carboxyle & acide \dots-oïque & acide éthanoïque, \ce{CH3-COOH}\\
ester & \chemfig[atom sep=1.6em]{-[:30](=[:90]O)-[:-30]O-[:30]} & \dots-oate de \dots-yle & éthanoate de méthyle, \ce{CH3-COO-CH3}\\
amine & \chemfig[atom sep=1.6em]{-[:30]N(-[:90]H)-[:-30]H} \;amino & \dots-amine & éthanamine, \ce{CH3-CH2-NH2}\\
amide & \chemfig[atom sep=1.6em]{-[:30](=[:90]O)-[:-30]N(-[:30]H)-[:-90]H} & \dots-amide & éthanamide, \ce{CH3-CO-NH2}\\
halogénoalcane & \chemfig[atom sep=1.6em]{-[:30]X} \;(X $=$ F, Cl, Br, I) & halogéno-\dots & chloroéthane, \ce{CH3-CH2-Cl}\\\bottomrule
\end{tabularx}}
\end{propriete}
\essSection{3}{La spectroscopie infrarouge}
\begin{propriete}[Bandes d'absorption caractéristiques]
Un spectre infrarouge représente la transmittance en fonction du nombre d'onde $\sigma$. Chaque type de liaison absorbe
dans un domaine caractéristique :
\par\smallskip
{\small\renewcommand{\arraystretch}{1.25}
\begin{tabularx}{\linewidth}{@{}lYl@{}}\toprule
\otTete{Liaison} & \otTete{Nombre d'onde (cm$^{-1}$)} & \otTete{Allure}\\\midrule
O--H (alcool) & 3200 à 3500 & bande forte et large\\
O--H (acide carboxylique) & 2500 à 3200 & bande très large\\
N--H (amine, amide) & 3100 à 3500 & bande moyenne\\
C--H & 2800 à 3100 & bandes fortes\\
C$=$O (aldéhyde, cétone, acide, ester, amide) & 1650 à 1750 & bande forte et fine\\\bottomrule
\end{tabularx}}
\end{propriete}
\schema{11-ir}
\begin{savoirfaire}[Exploiter un spectre infrarouge]
Chercher d'abord la bande fine et intense du groupe C$=$O vers \qty{1700}{cm^{-1}}, puis les bandes larges O--H ou
N--H au-dessus de \qty{2500}{cm^{-1}}. Croiser présences et absences pour proposer une famille, en cohérence avec la
formule brute.
\end{savoirfaire}
\essSection{4}{Les polymères}
\begin{definition}[Polymère et motif]
Un \textbf{polymère} est une macromolécule formée par la répétition d'un \textbf{motif}, issu d'une petite molécule
appelée monomère. Lors d'une polymérisation par addition, la double liaison C$=$C du monomère s'ouvre :
\[n\;\ce{CH2=CH2}\;\longrightarrow\;\ce{-(CH2-CH2)_{$n$}-}\qquad\text{(polyéthylène)}.\]
L'entier $n$, souvent de plusieurs milliers, est le \textbf{degré de polymérisation} : la masse molaire de la chaîne
vaut $n$ fois celle du motif.
\end{definition}
\begin{codecote}[0.56]
\begin{pythoncol}
M_motif = 28.0   # CH2-CH2 (g/mol)
def M_chaine(n):
    return n * M_motif
for n in (1, 1000, 40000):
    print(n, M_chaine(n))
\end{pythoncol}
\begin{consolecol}
1 28.0
1000 28000.0
40000 1120000.0
\end{consolecol}
\cote
\fontsize{9.5}{12.5}\selectfont
\capacite{Capacité numérique}\par\smallskip
La masse molaire d'une chaîne de polyéthylène est proportionnelle au degré de polymérisation $n$. Les chaînes d'un
sac plastique ($n\approx1000$) sont bien plus courtes que celles d'une prothèse de hanche en polyéthylène de très
haute masse molaire ($n\approx40\,000$) : les propriétés mécaniques en dépendent directement.
\end{codecote}
\begin{attention}
Ne pas confondre le monomère (\ce{CH2=CH2}, qui possède une double liaison) et le motif (\ce{-CH2-CH2-}, qui n'en a
plus). Pour retrouver le motif d'un polymère d'addition, repérer le plus petit fragment qui se répète le long de la
chaîne.
\end{attention}
\begin{rediger}
Nommer en détachant les indices par des tirets : « 3-méthylbutan-1-ol ». Pour un spectre infrarouge, citer la bande
(liaison, nombre d'onde, allure) avant de conclure : « bande fine et intense à \qty{1715}{cm^{-1}} : liaison C$=$O d'une
cétone, en l'absence de bande O--H ».
\end{rediger}
\exercices

\rubrique{Découverte}
\begin{blocLivre}{ex-11-1}
\begin{exercice}[1]{Reconnaître les familles fonctionnelles}
On considère les espèces suivantes.
\begin{center}\renewcommand{\arraystretch}{1.4}
\begin{tabular}{@{}l@{\hspace{10mm}}l@{\hspace{10mm}}l@{}}
\otLabel{A.}~\ce{CH3-CH2-CH2-OH} & \otLabel{B.}~\ce{CH3-CO-CH2-CH3} & \otLabel{C.}~\ce{HCOOH}\\
\otLabel{D.}~\ce{CH3-COO-CH2-CH3} & \otLabel{E.}~\ce{CH3-CH2-NH2} & \otLabel{F.}~\ce{CH3-CO-NH2}
\end{tabular}\end{center}
\begin{questions}
\item Recopier chaque formule, entourer le groupe caractéristique et nommer la famille correspondante.
\item Quelles espèces contiennent un groupe carbonyle C$=$O ?
\item Écrire la formule brute de l'espèce D. Quelle molécule d'une autre famille, déjà citée dans ce chapitre,
possède la même formule brute ?
\end{questions}
\end{exercice}
\end{blocLivre}
\begin{blocLivre}{ex-11-2}
\begin{exercice}[1]{Nommer des molécules}
Le document donne les formules topologiques de cinq molécules.
\schema{11-ex02}
\begin{questions}
\item Donner la famille de chacune des cinq molécules.
\item Les nommer.
\item Écrire la formule semi-développée des molécules A et D.
\end{questions}
\end{exercice}
\end{blocLivre}
\begin{blocLivre}{ex-11-3}
\begin{exercice}[1]{Du nom à la formule}
\begin{questions}
\item Écrire la formule semi-développée des molécules suivantes : 2-méthylbutan-1-ol ; pentan-3-one ; acide
butanoïque ; éthanamide ; 2-bromopropane.
\item Écrire la formule topologique du propanoate d'éthyle. Entourer la partie issue de l'acide et celle issue de
l'alcool.
\item Donner la formule brute de la pentan-3-one et celle du 2-méthylbutan-1-ol.
\end{questions}
\end{exercice}
\end{blocLivre}
\begin{blocLivre}{ex-11-4}
\begin{exercice}[1]{Des isomères de constitution}
Le document représente quatre molécules de formule brute \ce{C4H10O}.
\schema{11-ex04}
\begin{questions}
\item Vérifier que les quatre molécules ont bien la même formule brute. Comment les qualifie-t-on ?
\item Nommer les molécules A, B et C.
\item La molécule D appartient à une autre famille, celle des éthers. Qu'est-ce qui la distingue des trois autres ?
\item Citer une propriété qui peut différer d'un isomère à l'autre.
\end{questions}
\end{exercice}
\end{blocLivre}
\begin{blocLivre}{ex-11-5}
\begin{exercice}[1]{Lire un spectre infrarouge}
Le spectre infrarouge d'un liquide incolore est donné ci-dessous ; sa formule brute est \ce{C2H6O}.
\schema{11-ex05}
\begin{questions}
\item Attribuer les bandes repérées 1 et 2 à des liaisons de la molécule.
\item Quelle famille la bande 1 révèle-t-elle ? Identifier la molécule.
\item Le spectre de l'éthanal \ce{CH3-CHO} présenterait-il la bande 1 ? Quelle autre bande apparaîtrait ?
\end{questions}
\end{exercice}
\end{blocLivre}
\rubrique{Entraînement}
\begin{blocLivre}{ex-11-6}
\begin{exercice}[2]{Un arôme de poire}
L'arôme de poire des confiseries est dû à l'éthanoate de pentyle.
\begin{questions}
\item Écrire la formule semi-développée et la formule topologique de cette molécule, puis entourer son groupe
caractéristique. À quelle famille appartient-elle ?
\item Déterminer sa formule brute et calculer sa masse molaire.
\item Le propanoate de butyle a la même formule brute. Écrire sa formule semi-développée. Comment qualifier ces deux
molécules ?
\item L'éthanoate de pentyle peut être synthétisé à partir d'un acide carboxylique et d'un alcool (chapitre~7). Les
nommer.
\end{questions}
\begin{donnees}
$M(\ce{C})=\qty{12,0}{g.mol^{-1}}$ ; $M(\ce{H})=\qty{1,0}{g.mol^{-1}}$ ; $M(\ce{O})=\qty{16,0}{g.mol^{-1}}$.
\end{donnees}
\end{exercice}
\end{blocLivre}
\begin{blocLivre}{ex-11-7}
\begin{exercice}[2]{Trois flacons pour trois spectres}
Trois flacons étiquetés propan-1-ol, propanone et acide propanoïque ont perdu leurs étiquettes. Leurs spectres
infrarouges sont reproduits ci-dessous.
\schema{11-ex07}
\begin{questions}
\item Rappeler les bandes caractéristiques attendues pour chacune des trois familles.
\item Attribuer chaque spectre à une molécule, en justifiant par deux arguments à chaque fois.
\item Pourquoi la spectroscopie infrarouge ne permettrait-elle pas de distinguer le propan-1-ol du butan-1-ol ?
\end{questions}
\end{exercice}
\end{blocLivre}
\begin{blocLivre}{ex-11-8}
\begin{exercice}[2]{Des groupes dans les médicaments}
Le document donne les formules topologiques de trois médicaments courants.
\schema{11-ex08}
\begin{questions}
\item Recopier les formules, puis entourer et nommer les groupes caractéristiques du paracétamol.
\item Faire de même pour l'aspirine. Quelle famille, déjà rencontrée au chapitre~7, reconnaît-on ?
\item L'ibuprofène ne possède qu'un groupe caractéristique : lequel ? Quelle propriété acido-basique en découle
(chapitre~9) ?
\item Vérifier que la formule brute de l'ibuprofène est \ce{C13H18O2}.
\end{questions}
\end{exercice}
\end{blocLivre}
\begin{blocLivre}{ex-11-9}
\begin{exercice}[2]{La longueur des chaînes du polyéthylène}
Le polyéthylène des prothèses articulaires a des chaînes bien plus longues que celui des films plastiques. Le
programme recherche le degré de polymérisation d'une chaîne de masse molaire \qty{56000}{g.mol^{-1}}.
\begin{python}
M_motif = 28.0        # CH2-CH2 (g/mol)
n = 1
while n * M_motif < 56000:
    n = n + 1
print(n, n * M_motif)
\end{python}
\begin{console}
2000 56000.0
\end{console}
\begin{questions}
\item \capacite{Capacité numérique}Écrire l'équation de la polymérisation de l'éthène \ce{CH2=CH2} et préciser le
motif du polyéthylène.
\item Justifier la valeur 28,0 utilisée par le programme.
\item Interpréter l'affichage, puis retrouver le résultat par un calcul direct.
\item Combien d'atomes de carbone une telle chaîne contient-elle ?
\end{questions}
\end{exercice}
\end{blocLivre}
\begin{blocLivre}{ex-11-10}
\begin{exercice}[2]{Le PVC}
Le polychlorure de vinyle (PVC) des canalisations est obtenu par polymérisation du chloroéthène \ce{CH2=CHCl}.
\begin{donnees}
$M(\ce{C})=\qty{12,0}{g.mol^{-1}}$ ; $M(\ce{H})=\qty{1,0}{g.mol^{-1}}$ ; $M(\ce{Cl})=\qty{35,5}{g.mol^{-1}}$.
\end{donnees}
\begin{questions}
\item Écrire l'équation de la polymérisation et la formule du motif.
\item Calculer la masse molaire du motif, puis le pourcentage massique de chlore dans le PVC.
\item Une chaîne de PVC a une masse molaire de \qty{1,0e5}{g.mol^{-1}}. Calculer son degré de polymérisation.
\item L'incinération du PVC produit du chlorure d'hydrogène. Pourquoi les fumées doivent-elles être traitées ?
\end{questions}
\end{exercice}
\end{blocLivre}
\begin{blocLivre}{ex-11-11}
\begin{exercice}[2]{Un polyamide, le nylon}
Le nylon-6,6 est un polyamide : son motif, \ce{-CO-(CH2)4-CO-NH-(CH2)6-NH-}, relie les chaînes par des groupes
amide, comme les protéines.
\begin{donnees}
Formule brute du motif : \ce{C12H22N2O2} ; $M(\ce{C})=\qty{12,0}{g.mol^{-1}}$ ; $M(\ce{H})=\qty{1,0}{g.mol^{-1}}$ ;
$M(\ce{N})=\qty{14,0}{g.mol^{-1}}$ ; $M(\ce{O})=\qty{16,0}{g.mol^{-1}}$.
\end{donnees}
\begin{questions}
\item Recopier le motif et entourer les groupes amide.
\item Calculer la masse molaire du motif.
\item Une fibre de nylon a des chaînes de masse molaire voisine de \qty{4,5e4}{g.mol^{-1}}. Estimer leur degré de
polymérisation.
\item Les groupes amide créent des liaisons entre chaînes voisines. Quel intérêt pour une fibre textile ?
\end{questions}
\end{exercice}
\end{blocLivre}
\rubrique{Objectif bac}
\begin{blocLivre}{ex-11-12}
\begin{exercice}[3]{La carte d'identité de l'ibuprofène}
\begin{docu}{1}{Un anti-inflammatoire}
L'ibuprofène, de formule brute \ce{C13H18O2}, est le principe actif de nombreux anti-inflammatoires. Un comprimé
courant en contient 400~mg. Son spectre infrarouge présente une bande fine et intense à \qty{1710}{cm^{-1}} et une
bande très large entre 2500 et \qty{3200}{cm^{-1}}.
\end{docu}
\begin{donnees}
$M(\ce{C13H18O2})=\qty{206,3}{g.mol^{-1}}$ ; l'ibuprofène est un acide noté AH, de $\pKa=4{,}9$ ; solution titrante :
hydroxyde de sodium à $c_\mathrm{B}=\qty{0,100}{mol.L^{-1}}$.
\end{donnees}
\begin{questions}
\item Quelles informations les deux bandes du spectre infrarouge apportent-elles ? Quelle famille confirment-elles ?
\item Calculer la quantité de matière d'ibuprofène contenue dans un comprimé.
\item On dissout un comprimé et on titre l'ibuprofène par la soude, selon \ce{AH + HO- -> A- + H2O}. Calculer le
volume versé à l'équivalence.
\item \initiative Dans l'intestin, le pH vaut environ 6,5. Sous quelle forme l'ibuprofène s'y trouve-t-il
majoritairement ? Justifier.
\end{questions}
\end{exercice}
\end{blocLivre}
\begin{blocLivre}{ex-11-13}
\begin{exercice}[3]{Du PET à la fibre polaire}
\begin{docu}{1}{Un polyester recyclable}
Le poly(téréphtalate d'éthylène), ou PET, est le polymère des bouteilles transparentes. Recyclé, il est refondu puis
filé pour fabriquer des fibres textiles : une vingtaine de bouteilles suffisent pour une veste polaire. Son motif a
pour formule brute \ce{C10H8O4}.
\end{docu}
\schema{11-ex13}
\begin{donnees}
$M(\text{motif})=\qty{192,2}{g.mol^{-1}}$ ; masse de PET d'une bouteille : 25~g ; masse molaire moyenne des chaînes :
\qty{2,4e4}{g.mol^{-1}}.
\end{donnees}
\begin{questions}
\item Sur le motif représenté, repérer et nommer les groupes caractéristiques. Pourquoi parle-t-on de polyester ?
\item Calculer la quantité de motifs contenue dans une bouteille.
\item Calculer le degré de polymérisation moyen des chaînes.
\item \initiative Estimer le nombre de chaînes de polymère contenues dans une bouteille.
\item Citer un avantage du recyclage dit « chimique », qui redécoupe les chaînes en monomères, par rapport au
recyclage par simple refusion.
\end{questions}
\end{exercice}
\end{blocLivre}
\begin{blocLivre}{ex-11-14}
\begin{exercice}[3]{Les esters et les odeurs de fruits}
\begin{docu}{1}{Des arômes de synthèse}
De nombreux arômes alimentaires sont des esters :
\par\smallskip
{\small\renewcommand{\arraystretch}{1.25}
\begin{tabular}{@{}lll@{}}\toprule
\otTete{Ester} & \otTete{Odeur} & \otTete{Formule}\\\midrule
méthanoate d'éthyle & rhum & \ce{H-COO-C2H5}\\
butanoate d'éthyle & ananas & \ce{CH3-CH2-CH2-COO-C2H5}\\
éthanoate de 3-méthylbutyle & banane & \ce{CH3-COO-CH2-CH2-CH(CH3)2}\\\bottomrule
\end{tabular}}
\end{docu}
\begin{donnees}
$M(\ce{C})=\qty{12,0}{g.mol^{-1}}$ ; $M(\ce{H})=\qty{1,0}{g.mol^{-1}}$ ; $M(\ce{O})=\qty{16,0}{g.mol^{-1}}$.
\end{donnees}
\begin{questions}
\item Écrire la formule topologique du butanoate d'éthyle et entourer son groupe caractéristique.
\item Déterminer la formule brute du butanoate d'éthyle, puis celle de l'éthanoate de butyle. Comment qualifier ces
deux esters ?
\item De quel acide carboxylique et de quel alcool l'arôme de banane peut-il être obtenu ? Écrire l'équation de la
réaction d'estérification (chapitre~7).
\item \initiative Un flacon d'arôme contient 5,0~mL d'éthanoate de 3-méthylbutyle, de masse volumique
\qty{0,87}{g.mL^{-1}}. Calculer la quantité de matière correspondante.
\end{questions}
\end{exercice}
\end{blocLivre}
\begin{blocLivre}{ex-11-15}
\begin{exercice}[3]{Une molécule mystère}
\begin{docu}{1}{Les indices du laboratoire}
Un flacon sans étiquette contient un liquide incolore à l'odeur fruitée de dissolvant. La combustion complète de
0,881~g de ce liquide produit 1,76~g de dioxyde de carbone et 0,721~g d'eau ; le composé ne contient que du carbone,
de l'hydrogène et de l'oxygène. Une mesure donne une masse molaire de \qty{88}{g.mol^{-1}}. Le spectre infrarouge
montre une bande fine et intense à \qty{1740}{cm^{-1}} et aucune bande au-delà de \qty{3100}{cm^{-1}}.
\end{docu}
\begin{donnees}
$M(\ce{C})=\qty{12,0}{g.mol^{-1}}$ ; $M(\ce{H})=\qty{1,0}{g.mol^{-1}}$ ; $M(\ce{O})=\qty{16,0}{g.mol^{-1}}$ ;
$M(\ce{CO2})=\qty{44,0}{g.mol^{-1}}$ ; $M(\ce{H2O})=\qty{18,0}{g.mol^{-1}}$.
\end{donnees}
\begin{questions}
\item Calculer les quantités de matière de carbone et d'hydrogène contenues dans l'échantillon brûlé.
\item En déduire la masse d'oxygène de l'échantillon, puis la quantité de matière correspondante.
\item \initiative Déterminer la formule brute du composé.
\item Que révèle le spectre infrarouge ? Parmi les familles possibles, laquelle retenir ?
\item Proposer une formule semi-développée et un nom compatibles avec tous les indices, sachant que l'odeur évoque un
dissolvant pour vernis à ongles.
\end{questions}
\end{exercice}
\end{blocLivre}
